% Documento: DSSMV_ProjectDroid_report_1251486_1251506
% Relatório final do projeto PlantKeep (DSSMV, ISEP LETI)
% Turma 2DD, grupo #8 — Tiago Santos (1251486), Tomás Bastos (1251506)
% Template mandatória: SurveyFormat (llncs.cls + typeinst.tex)
\documentclass[runningheads,a4paper]{llncs}

\usepackage[utf8]{inputenc}
\usepackage[T1]{fontenc}
\usepackage{lmodern}
\usepackage{graphicx}
\usepackage{booktabs}
\usepackage{amsmath}
\usepackage{amssymb}
\usepackage{enumitem}
\setlist[itemize]{noitemsep, topsep=3pt, parsep=0pt, partopsep=0pt}

\renewcommand{\refname}{Bibliografia}

\begin{document}

\mainmatter

\title{PlantKeep\\ Relatório}

\titlerunning{PlantKeep}

\author{Tiago Santos (1251486) \and Tomás Bastos (1251506)}
\authorrunning{PlantKeep — Relatório}

\institute{ISEP — Licenciatura em Engenharia de Telecomunicações e Informática}

\maketitle

\begin{abstract}
Este relatório documenta o desenvolvimento do PlantKeep, uma aplicação Android em Java para a gestão de plantas de vaso e hortas urbanas. A aplicação combina dados das APIs públicas \textit{OpenPlantbook} e \textit{Open-Meteo}/\textit{IPMA} com características do dispositivo (sensor de luz ambiente, GPS e notificações) e uma base de dados remota \textit{restdb.io}, com cache local \textit{Room}. O documento apresenta a análise de requisitos (user stories, requisitos funcionais e não funcionais segundo FURPS+), o design de software em camadas (model/dto/handler/services) e a implementação em Java Android.
\\[2mm]
\textbf{Palavras-chave:} Android, Java, REST, OpenPlantbook, Open-Meteo, IPMA, Sensor de Luz, Arquitetura em Camadas, FURPS+.
\end{abstract}

% -------------------------------------------------------------
\section{Contexto}

O objetivo deste projeto é criar uma aplicação Android, para a gestão de plantas numa horta ou jardim. Quem cultiva plantas em casa encontra dificuldades: saber quando regar, saber se a temperatura ou a humidade do solo são adequadas à espécie, e acompanhar a evolução de cada planta ao longo do tempo. A aplicação centraliza toda esta informação e transformá-la em decisões práticas e automáticas.

O utilizador pode registar cada uma das suas plantas, indicando a espécie, e consultar a informação dessa espécie obtida da API OpenPlantbook  (ex: limites de temperatura e de humidade do solo recomendados, uma fotografia de referência). A aplicação combina esses limites com a previsão meteorológica para criar avisos: temperatura mínima ou máxima fora do intervalo previsto, humidade do solo abaixo do necessário (altura de regar) ou solo encharcado. Para cada planta é também determinada a próxima rega recomendada, tendo em conta a previsão de seca e o efeito de uma rega manual.

A aplicação utiliza características do telemóvel: o sensor de luz ambiente permite ao utilizador medir a luminosidade do local onde a planta está e verificar se é adequada à espécie. A localização (GPS) é usada para obter a previsão meteorológica correta para o localização.

A persistência está dividida em duas camadas: os dados do utilizador (plantas, registos do diário) são guardados na base de dados remota restdb.io, os dados meteorológicos e de espécie são externos, obtidos das APIs Open-Meteo e OpenPlantbook, respetivamente.

Este relatório documenta o contexto, a análise de requisitos, o design de software (arquitetura em camadas com separação model,dto,handler,service), a implementação em Java Android e os testes realizados.

% -------------------------------------------------------------
\section{Análise}

\subsection{User Stories}

\begin{itemize}
    \item \textbf{US01:} Como utilizador, eu quero registar uma nova planta (nome e espécie) para a poder acompanhar no futuro.
    \item \textbf{US02:} Como utilizador, eu quero consultar a lista das minhas plantas para ver rapidamente o estado de cada uma.
    \item \textbf{US03:} Como utilizador, eu quero consultar informação da espécie de uma planta (limites de temperatura e de humidade do solo, e uma imagem) para saber como a cuidar.
    \item \textbf{US04:} Como utilizador, eu quero saber a próxima rega recomendada de cada planta para saber quando a devo regar.
    \item \textbf{US05:} Como utilizador, eu quero registar manualmente uma rega para atualizar o histórico de cada planta.
    \item \textbf{US06:} Como utilizador, eu quero receber avisos quando a temperatura prevista sair do intervalo adequado à espécie para a proteger do frio ou do calor.
    \item \textbf{US07:} Como utilizador, eu quero receber avisos quando a humidade do solo prevista sair do intervalo adequado (seca ou solo encharcado) para saber se devo regar ou se o solo está saturado.
    \item \textbf{US08:} Como utilizador, eu quero medir a luminosidade ambiente do local onde está a planta e saber se é adequada à espécie para a colocar no sítio certo.
    \item \textbf{US09:} Como utilizador, eu quero manter um diário por planta, com observações e fotografias, para acompanhar a sua evolução ao longo do tempo.
    \item \textbf{US10:} Como utilizador, eu quero que a aplicação funcione com os dados em cache mesmo sem ligação à Internet.
\end{itemize}

\subsection{Requisitos Funcionais (RF)}

\begin{itemize}
    \item \textbf{RF01:} O sistema deve permitir criar uma nova planta com um nome (mínimo de 3 caracteres) e indicar a sua espécie.
    \item \textbf{RF02:} O sistema deve listar todas as plantas registadas pelo utilizador.
    \item \textbf{RF03:} O sistema deve obter da API pública \textit{OpenPlantbook} a informação da espécie (limites de temperatura e de humidade do solo, e imagem), identificada pelo seu \texttt{pid}.
    \item \textbf{RF04:} O sistema deve calcular a próxima rega recomendada para cada planta, com base na previsão meteorológica e no efeito de uma rega manual.
    \item \textbf{RF05:} O sistema deve registar a última rega manual de cada planta.
    \item \textbf{RF06:} O sistema deve produzir avisos de temperatura quando a previsão sair do intervalo [mín, máx] da espécie.
    \item \textbf{RF07:} O sistema deve produzir avisos de humidade do solo quando a previsão sair do intervalo [mín, máx] da espécie (seca ou solo encharcado).
    \item \textbf{RF08:} O sistema deve medir a luminosidade ambiente através do sensor de luz do dispositivo e avaliar a sua adequação à espécie.
    \item \textbf{RF09:} O sistema deve permitir adicionar registos ao diário de uma planta, incluindo fotografias.
    \item \textbf{RF10:} O sistema deve guardar os dados do utilizador na base de dados remota \textit{restdb.io}.
    \item \textbf{RF11:} O sistema deve manter uma cache local (\textit{Room}) para permitir o funcionamento offline.
    \item \textbf{RF12:} O sistema deve emitir notificações quando um aviso é disparado.
\end{itemize}

\subsection{Requisitos Não Funcionais (FURPS+)}

\subsubsection{Functionality (F)}
\begin{itemize}
    \item \textbf{RNF01 (Persistência):} Os dados do utilizador (plantas, registos do diário, limiares e preferências) devem ser guardados na base de dados \textit{restdb.io}, com uma cache local \textit{Room} como falha de segurança para funcionamento offline.
\end{itemize}

\subsubsection{Usability (U)}
\begin{itemize}
    \item \textbf{RNF02 (Interface):} A aplicação deve apresentar uma interface gráfica simples e intuitiva, em ecrãs dedicados (criar planta, lista, informação, rega, diário, luz), navegáveis a partir de um menu principal.
\end{itemize}

\subsubsection{Reliability (R)}
\begin{itemize}
    \item \textbf{RNF03 (Funcionamento offline):} A aplicação deve continuar a mostrar os dados previamente obtidos quando não há ligação à Internet, usando a cache local.
    \item \textbf{RNF04 (Robustez):} Uma falha de rede ou de API numa planta não deve bloquear o cálculo dos avisos das restantes plantas.
    \item \textbf{RNF05 (Validação):} O sistema deve validar os dados introduzidos pelo utilizador e recusar entradas inválidas (ex.: nome com menos de 3 caracteres).
\end{itemize}

\subsubsection{Performance (P)}
\begin{itemize}
    \item \textbf{RNF06 (Responsividade):} As operações de rede devem ser executadas em segundo plano (numa \texttt{Thread}), para que a interface responda de imediato às interações do utilizador.
\end{itemize}

\subsubsection{Supportability (S)}
\begin{itemize}
    \item \textbf{RNF07 (Portabilidade/Build):} O projeto deve ser compilável com \textit{Gradle} (AGP) num ambiente de desenvolvimento Android padrão.
\end{itemize}

\subsubsection{+ (Restrições de Implementação)}
\begin{itemize}
    \item \textbf{RNF08 (Linguagem):} A aplicação deve ser nativa para Android, escrita em \textbf{Java}.
    \item \textbf{RNF09 (Arquitetura):} O código deve seguir uma arquitetura em camadas, com separação \texttt{model}/\texttt{dto}/\texttt{handler}/\texttt{services}, em que o \texttt{model} é puro (sem rede, JSON ou dependências Android).
    \item \textbf{RNF10 (Comunicação com o exterior):} A aplicação deve comunicar com APIs externas via \textit{REST} (\textit{restdb.io}, \textit{OpenPlantbook}, \textit{Open-Meteo}/\textit{IPMA}) e explorar características do dispositivo (sensor de luz, GPS e notificações).
\end{itemize}

% -------------------------------------------------------------
\section{Design}


% -------------------------------------------------------------
\section{Implementação}


% -------------------------------------------------------------
\section{Testes}


\end{document}
